\section{Sistema de Code Review}

\subsection{Momento y principios del Review}

\begin{importantbox}{{Review Early, Review Often}}
\begin{itemize}
    \item No se revisa solo antes de fusionar con la rama principal, sino que \textbf{después de completar cada subtarea de grano fino} se debe activar la revisión de manera obligatoria.
    \item Evita que un error de diseño «contamine» la siguiente tarea.
    \item Entrada del Code Review: los requisitos/el plan original del usuario + el Git Diff.
\end{itemize}
\end{importantbox}

\subsection{Checklist para solicitar un Code Review}

\begin{enumerate}
    \item \textbf{Calidad del código}: separación de responsabilidades, desacoplamiento, manejo de errores, código duplicado, cobertura de casos límite.
    \item \textbf{Arquitectura}: si los roles de diseño son razonables, escalabilidad, impacto en el rendimiento, riesgos de seguridad.
    \item \textbf{Pruebas}: si se prueba la lógica real (y no solo con Mock), cobertura de casos límite, pruebas de integración.
    \item \textbf{Requisitos}: si se cumplen todos los requisitos del plan, sin proliferación de funcionalidades, y se alcanza un nivel de acabado apto para producción.
\end{enumerate}

\begin{warningbox}{Líneas rojas de conducta}
\begin{itemize}
    \item Está terminantemente prohibido decir «se ve bien» sin haber revisado el código.
    \item Está terminantemente prohibido marcar un problema secundario como un error fatal.
    \item Está terminantemente prohibido ser ambiguo — se debe llegar a una conclusión final.
\end{itemize}
\end{warningbox}

\subsection{Criterios para recibir un Code Review}

La esencia del Code Review es \textbf{una evaluación técnica basada en el código objetivo}, no una actuación social:

\begin{itemize}
    \item \textbf{Buscar el rigor técnico}: las decisiones se basan en hechos, resultados de pruebas y el estado actual de la arquitectura del sistema, no en la intuición o la autoridad.
    \item \textbf{Rechazar el acuerdo performativo}: no aceptar a ciegas las sugerencias del Reviewer sin haber verificado antes su viabilidad técnica.
    \item \textbf{Principio YAGNI}: si el Reviewer sugiere añadir una funcionalidad que el sistema actual no invoca en absoluto, debe rechazarse.
    \item \textbf{Verificar antes de actuar}: al recibir la retroalimentación, el primer paso es verificarla en la base de código, no modificar de inmediato.
    \item \textbf{Objeción técnica bien fundamentada} (Push Back): cuando se detecta que la sugerencia del Reviewer es errónea, se refuta con razonamiento técnico — la objeción no busca discutir, sino proteger la base de código.
\end{itemize}

\begin{knowledgebox}{Eliminar la ambigüedad}
Si el Reviewer da 6 observaciones y 2 de ellas no quedan claras, \textbf{está terminantemente prohibido hacer primero los 4 puntos que sí se entienden y dejar los 2 que no se entienden para el final}. Hay que detenerse de inmediato y preguntar primero todos los detalles con claridad, porque los módulos de software están altamente acoplados entre sí, y una comprensión parcial llevaría a una dirección de implementación equivocada, con un costo de retrabajo enorme.
\end{knowledgebox}

\begin{importantbox}{Los hechos hablan más que las palabras}
Cuando la observación del Reviewer es correcta, está terminantemente prohibido usar frases vacías de agradecimiento o elogio. Lo correcto es \textbf{corregirlo directamente y exponer los hechos de forma breve}. En una cultura de ingeniería eficiente, el código mismo es la mejor respuesta.
\end{importantbox}

\subsection{SOLID y la separación de responsabilidades}

Patrones de diseño arquitectónico que intervienen en el Code Review:

\begin{enumerate}
    \item \textbf{Principio de responsabilidad única} (SRP): evita que una clase o un archivo asuma todas las tareas; se debe desacoplar.
    \item \textbf{Principio de abierto/cerrado} (OCP): abierto a la extensión, cerrado a la modificación.
    \item \textbf{Principio de sustitución de Liskov} (LSP): una subclase puede sustituir a su clase padre.
    \item \textbf{Principio de segregación de interfaces} (ISP): las interfaces deben ser de grano fino.
    \item \textbf{Principio de inversión de dependencias} (DIP): depender de abstracciones y no de implementaciones concretas.
\end{enumerate}

La \textbf{separación de responsabilidades} es una filosofía arquitectónica de alto nivel: divide el software en módulos independientes con responsabilidades distintas, de modo que cada módulo resuelve solo un aspecto específico.

\subsection{Resumen de la sección}

El sistema de Code Review comprende dos Skills, «solicitar la revisión» y «recibir la revisión»; mediante el Checklist y las líneas rojas de conducta se garantiza el rigor técnico de la revisión, mientras que los principios SOLID y la separación de responsabilidades guían la evaluación de la arquitectura.

